\chapter{平台、地址空间与 I/O 端口}

\section{CPU 地址不等于 RAM}

物理地址空间由平台译码。某段地址可以落入 DRAM、ROM、PCI 窗口或空洞。
固件所在的高地址 ROM 和低地址 RAM 同时存在；把栈设为 \code{0x7000}
是平台布局决策，不是 x86 指令自动保证的事实。

\section{端口映射 I/O}

传统 PC 串口使用独立 I/O 端口空间，由 \code{in}、\code{out} 指令访问。
COM1 基址是 \code{0x3F8}。基址加偏移选择寄存器：数据寄存器偏移 0，
中断使能偏移 1，线路控制偏移 3，线路状态偏移 5。某些偏移的含义还会受
DLAB 位影响，因此寄存器访问是一段状态协议，不能看成互不相关的赋值。

\begin{contractbox}
软件写寄存器前必须知道设备状态；设备返回状态位后，软件才能推进所有权。
轮询必须有界，错误路径必须有可观察结果。
\end{contractbox}

\section{QEMU 在本项目中的职责}

QEMU 提供 \code{pc} 机器、\code{qemu64} CPU、TCG 执行、RAM、ISA 串口和
IDE 磁盘模型。它不替固件初始化串口，不替引导器读磁盘，也不替内核建立页表。
\code{-nodefaults} 减少无关默认设备，\code{-display none} 和
\code{-serial stdio} 把最早日志绑定到测试进程。
